\section{实战回顾与案例模板}
\subsection{代码修复示例}
Graham 描述了一个典型案例：Agent 发现 bug 之后，Plan 阶段写下要修改的文件列表，Act 阶段使用 \texttt{python apply\_patch} 生成 diff，Verify 阶段跑 \texttt{pytest}。全流程记录在日志里，若测试失败就自动回滚。

\begin{knowledgebox}{示例流程}
Plan：提取 Issue，列出文件与函数；Act：运行 Agent 写 patch，保存 diff；Verify：跑测试，若失败把 log 推送到 issue。
\end{knowledgebox}

\subsection{失败与纠错记录}
团队会把每次失败样本记录为“失败矩阵”：包含失败原因（测试、权限、语法）、Agent 输出、人工干预内容。该矩阵是训练 prompt 的重要素材。

\begin{warningbox}{失败日志不要删}
失败日志包含关键语义信息，用于生成 agent chains 训练数据；删除会导致未来无法复现故障。
\end{warningbox}

\subsection{本章小结}
\begin{itemize}
    \item 通过模版化流程可以把单次 success 定义为“Plan + Act + Verify” 三段完成。
    \item 记录失败矩阵为后续训练与 prompt 调整提供依据。
\end{itemize}

